\documentclass{article}

\usepackage{amsmath}
\usepackage{amssymb}
\usepackage{graphicx}

\title{\textbf{Latihan Soal Transformasi Geometri}}
\author{Toriq Afanudin}
\date{September 2026}

\begin{document}

\maketitle

\begin{enumerate}

    \item Bayangan dari titik $(2,3)$ oleh dilatasi dengan pusat $(4,1)$ dan faktor dilatasi $2$ adalah \dots

    \begin{enumerate}
        \item[A.] $(0,5)$
        \item[B.] $(5,0)$
        \item[C.] $(0,5)$
        \item[D.] $(5,5)$
        \item[E.] $(0,-5)$
    \end{enumerate}

    \item Bayangan dari titik $(2,3)$ oleh rotasi dengan pusat $(4,1)$ dan sudut putar $90^\circ$ berlawanan arah jarum jam adalah \dots

    \begin{enumerate}
        \item[A.] $(2,1)$
        \item[B.] $(-2,1)$
        \item[C.] $(2,-1)$
        \item[D.] $(1,2)$
        \item[E.] $(-1,2)$
    \end{enumerate}

    \item Titik $(4,-1)$ ditransformasikan oleh matriks
    $
        \begin{bmatrix}
            1 & 2 \\
            0 & 1
        \end{bmatrix}
    $,
    kemudian dilanjutkan dengan rotasi berpusat di $O(0,0)$ sebesar $90^\circ$ searah jarum jam.

    \begin{enumerate}
        \item Tentukan koordinat bayangan akhir titik tersebut.
        
        \item Tentukan persamaan garis yang memiliki gradien $2$ dan melalui titik bayangan akhir tersebut.
    \end{enumerate}
    \item Akan dicari bayangan dari $(-10,-3)$ yang dicerminkan terhadap sumbu x dilanjutkan dengan pencerminan terhadap sumbu y. Bayangannya adalah $\dots$

    \begin{enumerate}
        \item[A.] $(-3,-10)$
        \item[B.] $(-3,10)$
        \item[C.] $(3,-10)$
        \item[D.] $(3,10)$
        \item[E.] $(10,3)$
    \end{enumerate}

    \item Diketahui titik $A(3,-2)$ dipetakan oleh translasi $T=\begin{bmatrix}
        1\\-2
    \end{bmatrix}$ kemudian dilanjutkan oleh rotasi dengan pusat $(1,1)$ sejauh $90^{\circ}$. Koordinat titik hasil peta A adalah $\dots$

    \item Titik $(6,7)$ yang direfleksikan oleh titik $P(3,1)$ menghasilkan bayangan $\dots$

    \begin{enumerate}
        \item[A.] $(0,5)$
        \item[B.] $(5,0)$
        \item[C.] $(0,5)$
        \item[D.] $(5,5)$
        \item[E.] $(0,-5)$
    \end{enumerate}

    \item Komposisi dari 2 buah translasi adalah translasi.
    \begin{enumerate}
        \item[A.] Benar
        \item[B.] Salah
    \end{enumerate}
    \item Suatu titik A dengan koordinat $(2,3)$ jika di transformasikan oleh matriks $\begin{bmatrix}
        4&1\\1&4
    \end{bmatrix}$ kemudian dilanjutkan dengan pencerminan terhadap garis $y=x$ akan memiliki bayangan $\dots$
    \begin{enumerate}
        \item[A.] $(16,12)$
        \item[B.] $(14,11)$
        \item[C.] $(11,14)$
        \item[D.] $(9,11)$
        \item[E.] $(11,9)$
    \end{enumerate}

    \item Suatu titik $(x,y)$ ditransformasikan terhadap matriks $\begin{bmatrix}
        k&0\\0&k
    \end{bmatrix}$ penyelesaianya sama dengan perkalian skalar $k$ dengan $\begin{bmatrix}
        x\\y
    \end{bmatrix}$.
    \begin{enumerate}
        \item[A.] Benar
        \item[B.] Salah
    \end{enumerate}
    \item Transformasi rotasi terhadap titik asal $(0,0)$ sebesar $90^{\circ}$ berhubungan dengan matriks $\begin{bmatrix}
        0&1\\-1&0
    \end{bmatrix}$.
    \begin{enumerate}
        \item[A.] Benar
        \item[B.] Salah
    \end{enumerate}
    \item Semua bayangan hasil transformasi geometri bangun datar memiliki ukuran sama dengan bangun datar sebelumnya, kecuali translasi.
    \begin{enumerate}
        \item[A.] Benar
        \item[B.] Salah
    \end{enumerate}

\end{enumerate}

\end{document}